\subsection*{3.4. 汇总}

命题 3.1.2、3.2.1、3.2.2、3.2.3 及 3.3.5 构成“Picard--Lefschetz 公式”. 此公式为 XIV, 3.2.11 的代数类似物, 但其形式随相对维数的奇偶性而大不相同.

在最重要的 $X/S$ 固有的情形中, 结果可汇总如下.

\subsubsection*{定理 3.4 (Picard--Lefschetz 公式).}
\underline{设 $S$ 为严格 Hensel 特质, $f:X\to S$ 为纯相对维数 $n=2m$ 或 $2m+1$ 的固有平坦态射. 假设 $X_\eta$ 光滑, $X_s$ 仅有普通二次奇点. 设 $\Sigma$ 为 $X_s$ 的奇点集.}

\begin{enumerate}
\item[(i)] \underline{对每个 $x\in\Sigma$ 有一个消灭圈}
\[
\pm\delta_x\in H_n(X_{\bar\eta})(m).
\]
\underline{各 $\delta_x$ 两两正交. 当 $n$ 为偶数时,}
\[
(\delta_x,\delta_x)=(-1)^m,2.
\]

\item[(ii)] \underline{对 $i\ne n,n+1$, 有同构}
\[
H^i(X_s)\xrightarrow{\ \sim\ }H^i(X_{\bar\eta})
\]
\underline{并有正合列}
\[
0\longrightarrow H^n(X_s)
 \longrightarrow H^n(X_{\bar\eta})
 \xrightarrow{\,((\,\cdot\,,\delta_x))_{x\in\Sigma}\,}
 \bigoplus_{x\in\Sigma}\Lambda(n-m)
 \longrightarrow H^{n+1}(X_s)
 \longrightarrow H^{n+1}(X_{\bar\eta})
 \longrightarrow0.
\]

\item[(iii)] \underline{有惯性群 $I$ 的局部特征标:}
\[
\begin{array}{lll}
n\text{ 为偶数}:&\varepsilon_x:I\longrightarrow\{\pm1\},\\[2pt]
n\text{ 为奇数}:&\varepsilon_x:I\longrightarrow\Lambda(1).
\end{array}
\]
\underline{当 $X$ 在 $x$ 处正则时, 这些特征标是满射; 当 $x$ 非退化时, 它们分别为唯一的 $2$ 阶特征标与典范特征标.}

\underline{惯性群 $I$ 的作用为}
\[
n=2m\qquad
\sigma(a)=a+(-1)^m\sum_{x\in\Sigma}
 \frac{\varepsilon_x(\sigma)-1}{2}
 (a,\delta_x)\delta_x,
\]
\[
n=2m+1\qquad
\sigma(a)=a-(-1)^m\sum_{x\in\Sigma}
 \varepsilon_x(\sigma)(a,\delta_x)\delta_x.
\]
\underline{特别地,}
\[
\begin{array}{lll}
n\text{ 为偶数}:&\sigma(\delta_x)=\varepsilon_x(\sigma)\delta_x,\\[2pt]
n\text{ 为奇数}:&\sigma(\delta_x)=\delta_x.
\end{array}
\]
\end{enumerate}